\subsection{Typy statických analyzátorů}
\label{subsec:static-types}

Pod pojmem statický analyzátor se skrývá široká škála nástrojů, které se liší hloubkou analýzy a rozsahem zachycených problémů.
Pro účely srovnání s LLM postačí rozlišit tři základní úrovně, od nejjednodušších kontrol až po rozsáhlé enterprise systémy.

\begin{itemize}
    \item \textbf{Kompilátorová analýza} je nejzákladnější formou statické analýzy.
    Kompilátory jako GCC nebo Clang\footnote{\url{https://clang.llvm.org/}} při překladu upozorňují na zjevné problémy, například nevyužité či neinicializované proměnné.

    \item \textbf{Lintery a pokročilé analyzátory} rozšiřují tuto vrstvu o sady pravidel pro dobré programátorské zvyky, bezpečnost a styl.
    Typickými zástupci pro jazyk C++ jsou \texttt{clang-tidy}\footnote{\url{https://clang.llvm.org/extra/clang-tidy/}} a \texttt{cppcheck}\footnote{\url{https://cppcheck.sourceforge.io/}}.

    \item \textbf{Enterprise nástroje}, například SonarQube\footnote{\url{https://www.sonarsource.com/products/sonarqube/}}, propojují statickou analýzu s dlouhodobým sledováním kvality kódu.
    Pro akademické prostředí systému Kelvin jsou však předimenzované.
\end{itemize}

Ze všech uvedených kategorií jsou pro školní systém relevantní především lintery, protože nabízejí vyváženou kombinaci pokrytí a provozní jednoduchosti.
Samotné rozdíly mezi těmito přístupy a velkými jazykovými modely jsou podrobněji rozebrány v následující sekci.

\endinput
